\section{NREM and Memory Consolidation}\label{sec:consolidation}

While N3 sleep is cleaning the brain, it is also actively filing away important memories. NREM sleep is fundamental for \textit{consolidation}, the process of transforming new, fragile memories (typically held in the hippocampus) into stable, long-term memories (stored in the neocortex) \citep{rasch2007}.

This is achieved through a precise neural dialogue:
\begin{enumerate}
    \item \textbf{Hippocampal Replay:} The hippocampus, which recorded episodic events from the day (e.g.,~the path through a maze, a list of new vocabulary), replays these neural sequences at a highly compressed speed (up to 20x faster) \citep{wilson1994reactivation, ji2007}.
    \item \textbf{Spindles and Slow Waves:} These replays occur in coordination with the N2 sleep spindles and N3 slow-oscillations. This ``replay-and-spindle'' activity is believed to be the mechanism that ``teaches'' the neocortex, integrating the new information into long-term knowledge networks \citep{rasch2007}.
\end{enumerate}
Put simply, NREM sleep is like saving your work. It takes the ``in-progress'' files from the day and integrates them into the brain's permanent hard drive.
